% -----------------------------------------------------------------------------------
\section{Problems as programs}\label{sec:programs}

Every result so far has been proved on problems that \emph{supply} their counterfactuals. A problem
of Section~\ref{sec:model} comes with the law $\mu_\pi$ for every policy $\pi$, and with structural
plans that fix what each response would be on every branch. The decision theories that the debate
compares disagree precisely about where such laws come from: causal, evidential, and functional
theories differ over which parts of the world would have been different had the agent chosen
otherwise. A model that hands this dependence to the agent has already settled what the theories
dispute. This is why ex-ante optimality came so cheaply in Theorem~\ref{thm:universal}, and why the
deeper results of Part~\ref{part:across} concerned the shape of an evaluation rule rather than the
identification of alternatives.

This part removes the supply. An agent is now a computable procedure that receives a description of
its situation---by default the source code of the environment---and must itself determine what
would happen under each of its possible outputs. The question is whether some such procedure is
optimal on every fair problem. Answering it requires saying what \enquote{optimal} means once the
comparison is no longer given by a supplied law, and that is where most of the work lies.

\subsection{Agents and environments}

Fix an effective language of \emph{descriptions}, finite strings from which a finite nonempty set
$A(x)$ of legal actions can be computed for each description $x$. For the negative results a fixed
binary action set suffices, so nothing depends on how defaults are chosen when action sets vary.

\begin{definition}[Agents]\label{def:agent}
An \emph{agent} is a partial computable function $T$ from descriptions to actions with
$T(x)\in A(x)$ whenever $T(x)$ is defined. It is \emph{globally total} if it halts on every
well-formed description, and \emph{promise-total} on a class $\cls$ of environments if it halts on
every description belonging to $\cls$.
\end{definition}

The two notions differ. \enquote{A total computable function} normally means the first;
\enquote{an agent that halts on every problem in the class} permits the second. Both receive
answers below, and the distinction turns out to separate a positive theorem from a negative one.

\begin{definition}[Program environment]\label{def:progenv}
A \emph{program environment} $E$ is a program with a finite set $H_E$ of information sets, each
$h\in H_E$ carrying a finite nonempty action set $A_{E,h}$ and an \emph{owned query} $q(E,h)$, a
description from which $E$ and $h$ can be recovered. Its policy space is
$\Pol_E=\prod_{h\in H_E}A_{E,h}$, and its root description $d(E)$ is one of its owned queries. Given
access to an agent, $E$ produces a payoff. For a table $\pi\in\Pol_E$, $E^\pi$ denotes the execution
in which every owned query is answered by the corresponding entry of $\pi$.
\end{definition}

Four stipulations make owned queries behave like the information sets of
Definition~\ref{def:problem}: a query identifies its parent environment and its information set;
distinct environments have disjoint owned-query namespaces; repeated queries receive the same
answer; and, when the fairness condition $\cls_{\mathrm{self}}$ below is imposed, the environment
consults the agent only through its owned queries. The disjointness matters: if two environments
may both call one query their own, the conflict of Proposition~\ref{prop:foreign} below reappears
inside a class that is ostensibly self-only.

Three representations of payoffs must be kept apart. An execution may return an \emph{exact
rational}, or an explicit finite distribution with rational probabilities and payoffs; it may return
a \emph{computable real}, given by a program producing rational approximations of any required
accuracy; or it may be a random execution returning a \emph{sampled} payoff whose expectation is
the value. The positive results use the first; the strongest negative result
(Theorem~\ref{thm:exactsel}) uses only sampled payoffs in $\{0,1\}$, so it does not exploit exotic
real-valued outputs.

\subsection{Replacement and anchored evaluation}

To ask whether $T$ is optimal on $E$ one must say what $E$ would pay another agent $A$. Some ways of
saying it define a single payoff function $U(E,A)$ common to all agents. Others first run or inspect
a reference agent $T$, freeze part of the resulting situation, and only then consider alternatives;
these define an \emph{anchored kernel} $U_K(E,T;A)$, where $K$ names the convention. Every kernel is
required to satisfy the baseline identity
\begin{equation}\label{eq:baseline}
  U_K(E,T;T)=U(E,T),
\end{equation}
and the corresponding notion of optimality is
\begin{equation}\label{eq:anchored}
  U(E,T)\ge U_K(E,T;A)\qquad\text{for every admissible agent }A.
\end{equation}
In general there is no function $W$ with $U_K(E,T;A)=W(E,A)$ for every reference $T$. Some versions
of the question therefore ask for \emph{the best agent in a common environment}, others for \emph{an
agent that is a fixed point of its own counterfactual evaluations}. The difference is the
program-level form of the difference between \cri{1} and \cri{3}.

\begin{definition}[Replacement conventions]\label{def:conventions}
Let $E$ be an environment with a designated \emph{action port}, the place at which the agent's
actual action on $d(E)$ is received.
\begin{description}[font=\normalfont,leftmargin=2.4em,labelwidth=1.6em,labelsep=.8em,itemsep=.35em]
  \item[$O_{\mathrm{const}}$] Replace the action at the port by $A$'s action, and freeze every other
        response of the environment at its value under the reference $T$; fixed embedded algorithms
        are left unchanged.
  \item[$O_{\mathrm{oracle}}$] Replace every designated oracle call, the port included, by a call to
        $A$.
  \item[$O_{\mathrm{syn}}$] Replace specified, typed occurrences of program components whose syntax
        equals the reference program.
  \item[$O_{\mathrm{ext}}$] Replace specified program components that pass a stated test of
        extensional equivalence with the reference agent (Section~\ref{sec:extensional}).
\end{description}
Under $O_{\mathrm{syn}}$ and $O_{\mathrm{ext}}$ the description $d(E)$ supplied to the agent is held
fixed during replacement.
\end{definition}

The last clause is a choice. Supplying the rewritten source $d(E[A])$ instead would change the
agent's observation as well as the program occupying the agent's role; that is a legitimate model,
but a different one (Proposition~\ref{prop:rawsource} shows one of its consequences). Raw
substring replacement is not an adequate reading of $O_{\mathrm{syn}}$, since it can alter quoted
data, comments, or unrelated identifiers; typed occurrences or explicit component interfaces are
needed. Logical counterfactuals, $O_{\mathrm{log}}$, form a family of semantics rather than a fifth
replacement operation, and are treated in Section~\ref{sec:logical}.

\subsection{Fairness}

The fairness condition of Definition~\ref{def:fair} splits into several conditions, which overlap
rather than partition the environments.
\begin{description}[font=\normalfont,leftmargin=2.4em,labelwidth=1.6em,labelsep=.8em,itemsep=.35em]
  \item[$\cls_{\mathrm{self}}$] Dependence on the evaluated agent is confined to its answers on the
        owned queries $q(E,h)$.
  \item[$\cls_{\mathrm{any}}$] The environment may consult the agent on arbitrary descriptions.
  \item[$\cls_{\mathrm{ext}}$] Replacing agents by extensionally equal agents preserves the payoff
        kernel: $T\equiv T'$ and $A\equiv A'$ imply $U_K(E,T;A)=U_K(E,T';A')$.
  \item[$\cls_{\mathrm{const}}$] Fixed algorithms, including copies of a particular agent, may occur
        as constants of the environment.
\end{description}
Two locality requirements should also be separated. \emph{Executed-query locality} says that actual
executions call agents only on owned queries. \emph{Kernel locality} says that the whole
counterfactual kernel, including the rule that decides what to replace, depends on the reference
agent only through its answers on owned queries. Global program-equivalence matching
(Section~\ref{sec:extensional}) satisfies the first while violating the second.

In Section~\ref{sec:model} a fixed algorithm inside a problem was harmless: its output on the
relevant input is a number that is part of the problem, and the problem is a finite object
specified by such numbers. In a program environment the algorithm's code is part of the description
and its output is not; the agent must compute it, and may be the very thing being computed. This is
what makes $\cls_{\mathrm{const}}$ consequential here (Theorem~\ref{thm:fixedcopy}).

\subsection{Programs and causal matrices}

The conventions $O_{\mathrm{oracle}}$ and $O_{\mathrm{const}}$ have exact counterparts in the
causal matrix of Section~\ref{sec:matrix}. Let $E$ be a finite oracle environment in which every
execution with table answers halts with an exact rational payoff, and let the owned queries be
answered at two kinds of position: the action ports, and the remaining oracle calls. For
$\sigma,\tau\in\Pol_E$ let $M_E[\sigma,\tau]$ be the payoff of the execution in which the oracle
calls are answered by $\sigma$ and the ports by $\tau$.

\begin{proposition}[Programs induce causal matrices]\label{prop:progmatrix}
For agents $T$ and $A$ write $\pi_T$ and $\pi_A$ for the tables of their answers on the owned
queries. Then
\[
  U_{\mathrm{oracle}}(E,A)=M_E[\pi_A,\pi_A],\qquad
  U_{\mathrm{const}}(E,T;A)=M_E[\pi_T,\pi_A].
\]
$M_E$ is the causal matrix of the finite problem whose structural states are the oracle tables,
with $\mu_\pi=\delta_\pi$ and $U(\sigma,\tau)=M_E[\sigma,\tau]$, and every finite real square matrix
indexed by a product of finite action sets arises in this way.
\end{proposition}

\begin{proof}
Under $O_{\mathrm{oracle}}$ every owned query, port or not, is answered by $A$, which is the
execution indexed by $(\pi_A,\pi_A)$. Under $O_{\mathrm{const}}$ the ports receive $A$'s answers
while every other response is frozen at its value under $T$; since the only responses depending on
the agent are the oracle calls, this is the execution indexed by $(\pi_T,\pi_A)$. The problem
described has causal matrix $\sum_\rho\mu_\sigma(\rho)U(\rho,\tau)=M_E[\sigma,\tau]$. For the last
claim, an environment that calls the oracle on all owned queries, then reads the ports, and pays a
prescribed entry $M[\sigma,\tau]$ of a given matrix realises $M$, exactly as in
Theorem~\ref{thm:realise}.
\end{proof}

Oracle substitution thus recovers the diagonal of the matrix, ex-ante value, and port-only
intervention recovers a row, causal evaluation at the baseline $\pi_T$. Both require a designation
of which positions are ports and which are calls; a program in which the agent's code occurs inside
other programs, with no designated interface, does not carry that designation, and the agent must
construct it. Section~\ref{sec:operators} shows that the two kinds of evaluation reappear, in this
form, as the only two coherent ways of doing so.

\subsection{Comparing the conventions}

Replacing more components does not make optimality logically stronger or weaker.

\begin{proposition}[No implications between conventions]\label{prop:nonimplication}
For any two of $O_{\mathrm{const}}$, $O_{\mathrm{oracle}}$, $O_{\mathrm{syn}}$, $O_{\mathrm{ext}}$ there
are terminating environments on which a given agent is optimal under the first and not under the
second, and environments on which the reverse holds.
\end{proposition}

\begin{proof}
Let a typed environment contain an action port, an additional oracle call, a literal copy of the
reference agent's code, and an independently encoded copy computing the same function, each
queried at $d(E)$; let the reference agent be the constant $0$ and the alternative the constant $1$.
The four conventions then produce the following patterns of the four values.
\begin{center}
\small
\begin{tabular}{@{}lcccc@{}}
  \toprule
  evaluation & port & oracle call & literal copy & encoded copy \\
  \midrule
  baseline                & 0 & 0 & 0 & 0 \\
  $O_{\mathrm{const}}$ alternative  & 1 & 0 & 0 & 0 \\
  $O_{\mathrm{oracle}}$ alternative & 1 & 1 & 0 & 0 \\
  $O_{\mathrm{syn}}$ alternative    & 1 & 1 & 1 & 0 \\
  $O_{\mathrm{ext}}$ alternative    & 1 & 1 & 1 & 1 \\
  \bottomrule
\end{tabular}
\end{center}
For two conventions producing distinct patterns $v$ and $w$ from the baseline pattern $z$, set the
payoff to $1$ at $z$, $0$ at $v$, $2$ at $w$, and $0$ elsewhere. The reference agent is optimal under
the first convention and not under the second; exchanging the payoffs at $v$ and $w$ gives the
reverse.
\end{proof}

The conventions coincide only under conditions that make their replacement sets coincide.
Oracle and syntactic substitution agree under a slot-preserving compilation of oracle calls into
designated code occurrences (Section~\ref{sec:syntactic}); oracle and extensional substitution agree
when the extensional test selects exactly the designated dependencies; syntactic and extensional
substitution agree when their replacement masks coincide. For logical semantics no comparison is
possible until a particular semantics is fixed: a guarded provability semantics is partial
(Definition~\ref{def:guarded}), a proof of the value of each policy-substituted execution simply
evaluates an independently supplied substitution (Proposition~\ref{prop:proofopt}), and the scores
of a logical inductor are finite-time estimates. Treating these as one fifth payoff function would
conceal the main issue.
